% Generated draft: requires mathematical and source review.
\documentclass[11pt]{article}
\usepackage[margin=25mm]{geometry}
\usepackage{fontspec}
\setmainfont{Latin Modern Roman}
\usepackage{amsmath,amssymb,mathtools,unicode-math}
\setmathfont{Latin Modern Math}
\usepackage{enumitem,needspace,xurl,hyperref}
\setlength{\parindent}{0pt}
\setlength{\parskip}{5pt}
\setlength{\emergencystretch}{3em}
\newcommand{\N}{\mathbb{N}}
\newcommand{\Z}{\mathbb{Z}}
\newcommand{\Q}{\mathbb{Q}}
\newcommand{\R}{\mathbb{R}}
\newcommand{\C}{\mathbb{C}}
\newcommand{\F}{\mathbb{F}}
\newcommand{\A}{\mathbb{A}}
\newcommand{\PP}{\mathbb{P}}
\newcommand{\E}{\mathbb{E}}
\newcommand{\eps}{\varepsilon}
\newcommand{\dd}{\,\mathrm{d}}
\begin{document}
\section*{UnsolvedMath ID: 30003939}
\textbf{Source problem code:} catalogue entry
\subsection{Definitions and mathematical statement}
The supplied record concerns parameters $\beta>1$ and a functional denoted by $L_{\beta}$, described there as a candidate Lyapunov functional. It supplies neither a formula for $L_{\beta}$ nor the evolution equations whose long-time numerical solutions exhibit the structures in question.

For clarity, if an admissible topological space $X_{\beta}$ and a functional $L_{\beta}\colon X_{\beta}\to\R\cup\{+\infty\}$ were specified, a finite-energy local minimizer would mean a point $u_*\in X_{\beta}$ with $L_{\beta}(u_*)<+\infty$ for which there exists a neighborhood $U$ of $u_*$ such that
\[
L_{\beta}(u_*)\leq L_{\beta}(v)\qquad\text{for every }v\in U.
\]
This is an explanatory definition only: the space, topology, admissible constraints, and even the intended domain and codomain of $L_{\beta}$ are not specified by the supplied record.

\textbf{Source question.} For $\beta>1$, determine whether the singular structures observed in long-time numerical solutions are related to local minima of the candidate Lyapunov functional $L_{\beta}$.

Here ``singular structures,'' ``long-time,'' and ``related to'' have no mathematical definitions in the supplied material. In particular, the record does not specify whether the intended relation is equality with a minimizer, convergence to a minimizer, correspondence with the support of a minimizer, or another relation. Nor does it specify whether every observed structure, some structures, or all possible limiting structures are intended. Consequently, the source question cannot yet be expressed as a uniquely determined quantified proposition. The local-minimum definition above does not resolve these missing qualifications.
\subsection{Short English statement}
For $\beta>1$, the source asks how numerically observed long-time singular structures are connected with local minima of $L_{\beta}$. A precise theorem requires the functional, evolution, admissible states, topology, boundary conditions, and intended correspondence, none of which are supplied. The recorded partially solved status is supplied dated curation and requires independent review.
\subsection{Sources}
\begin{itemize}
\item Local registry: \texttt{data/unsolvedmath/problems.json}, permanent UnsolvedMath ID \texttt{30003939}, source problem code \texttt{catalogue entry}, title ``Singular Branched-Transport Structures as Lyapunov Minima.'' The record attributes the problem to Oberwolfach workshop open-problem discussions, gives proposed year 2018, and names no individual proposer.
\item Supplied report citation: \emph{Reactive Flows in Deformable, Complex Media} (2019), Oberwolfach Reports, pp.\ 2385--2473, DOI \texttt{10.4171/owr/2018/39}. Also listed as ``Oberwolfach Report 39/2018.'' \url{https://doi.org/10.4171/owr/2018/39}.
\item Supplied literature reference: \emph{Branching structures emerging from a continuous optimal transport model}, described in the record as a 2021 primary study. \url{https://doi.org/10.1016/j.jcp.2021.110700}.
\item Literature assessment checked at \texttt{2026-08-21}: this is supplied dated curation, not independently verified. It labels the problem partially solved and reports that the 2021 study derives a Lyapunov functional and numerically finds singular and fractal-like equilibria associated with its minima, while a general theorem identifying all such structures with local minima is not established. The supplied background calls this dated literature triage derived from OWR curation and requests independent verification of time-sensitive status and cited claims.
\end{itemize}
\end{document}
